% ===== tr-07 块6：第 5 章 pp.20--23（原文页 20--23）=====

% [原文 p.20]
\begin{itemize}
\item Calabi--Yau 体积测度与 Fubini--Study 度量的体积形式非常不同（参看第 5.3.1 节），因此“相对于 Fubini--Study 型环境度量的体积密度”这一想法不再合适。(5.10) 的替代品原来是 Skoda 型估计 (5.11)。
\item 与 $L^\infty$-界不同，Hölder 范数强烈依赖于环境度量的选取，而环境度量指定了距离函数的选取。我们认为如何在该族中一致地作出半显式的选取是高度不显然的，因此我们不试图推广 Hölder 估计。
\item 在我们的应用中有一个技术问题：涉及比较两个势时，只有一个势具有体积密度控制。因此与上面的 $L^1$-稳定性估计不同，我们的版本不对称地处理这两个势。
\end{itemize}

我们对势估计的类比是

\begin{theorem}{5.4.5}
（参看 [52, 第 2.2 节]）设 $(Y,\omega)$ 是紧 Kähler 流形，$\phi\in PSH(Y,\omega)\cap C^0$，且 $\omega_\phi^n$ 是绝对连续测度。假设存在正数常数 $\alpha$、$A$，使得关于 $\omega_\phi^n$ 的 Skoda 型估计成立：
\begin{equation}
\int_Y e^{-\alpha u}\frac{\omega_\phi^n}{\mathrm{Vol}(Y)}\le A,\qquad \forall u\in PSH(Y,\omega)\ \text{且}\ \sup_Y u=0. \tag{5.11}
\end{equation}
则我们有
\begin{itemize}
\item 若 $\sup_Y\phi=0$，则 $\|\phi\|_{C^0}\le C(n,\alpha,A)$。
\item 对固定的 $n,\alpha,A$，存在数 $B(n,\alpha,A)$，使得若对某个 $t_0$ 有
\[
\frac{\int_{\phi\le-t_0}\omega_\phi^n}{\mathrm{Vol}(Y)}<(2B)^{-2n},
\]
则
\[
\min\phi\ge -t_0-4B\Big(\frac{\int_{\phi\le-t_0}\omega_\phi^n}{\mathrm{Vol}(Y)}\Big)^{1/2n}.
\]
\end{itemize}
\end{theorem}

% [Tutor Note: (5.11) 的测度是 $\omega_\phi^n$（而不是 $\omega^n$），并且除以 $\mathrm{Vol}(Y)=\int_Y\omega^n$ 归一化；这与 (5.10) 用 $\omega^n$ 归一化不同。定理 5.4.5 的结论第二项的指数是 $1/2n$（不是 $1/n$ 或 $1/2$），且系数 $4B$ 与 $(2B)^{-2n}$ 中的 $2$ 都不能漏。]

这一结果的力量在于：即使复结构/Kähler 结构高度退化，它仍然适用，因为它把对 $(Y,\omega)$ 的依赖蒸馏为仅 3 个常数 $n,\alpha,A$。它与上面更标准版本的关系可解释如下：
\begin{itemize}
\item 对固定的环境 Kähler 度量，由 Hölder 不等式与标准 Skoda 不等式（定理 5.4.3），在密度可积性假设 (5.10) 之下，取 $\frac{1}{p}+\frac{1}{p'}=1$ 以及某个足够小的 $\beta>0$，有
\[
\resizebox{0.97\textwidth}{!}{$\displaystyle \int_Y e^{-\beta u}\frac{\omega_\phi^n}{\mathrm{Vol}(Y)} \le \frac{1}{\mathrm{Vol}(Y)}\Big(\int_Y\Big|\frac{\omega_\phi^n}{\omega^n}\Big|^p\omega^n\Big)^{1/p}\Big(\int_Y e^{-p'\beta u}\omega^n\Big)^{1/p'}\le \mathrm{const}.$}
\]
这意味着即使在标准情形下，Skoda 型估计 (5.11) 也是比 (5.10) 更弱的假设。
\end{itemize}

% [Tutor Note: 上式中的 Hölder 共轭指数是 $p$ 与 $p'$（$\frac1p+\frac1{p'}=1$），右端第二个因子里的指数是 $p'\beta u$ 而不是 $p\beta u$；第一项把 $e^{-\beta u}$ 与密度比配对，第二项用 Skoda 不等式吸收 $e^{-p'\beta u}$。]

% [原文 p.21]
\begin{itemize}
\item 结论的第一部分重新得到 Kolodziej 的势估计。
\item 结论的第二部分是关于比较两个势 $\phi$ 与 $-t_0$ 的。条件 $\frac{\int_{\phi\le-t_0}\omega_\phi^n}{\mathrm{Vol}(Y)}<(2B)^{-2n}$ 意味着 $\int_{\phi\le-t_0}\omega_\phi^n$ 很小；由 Hölder 不等式与密度可积性 (5.10)，这又可由 $\int_{\phi\le-t_0}\omega^n$ 的小性推出。这应看作 $L^1$-势稳定性条件 $\int_Y|\phi+t_0|\omega^n\ll1$ 的一半。关于 $\min\phi$ 下界的结论，应看作对 $\phi+t_0$ 的 $C^0$-范数的小性界的一半，即两个势 $\phi$ 与 $-t_0$ 在 $C^0$-范数下彼此接近。
\item 要推广到 $\psi$ 不必为常数的情形，只需在定理 5.4.5 中把 $\omega$ 换成 $\omega_\psi$、把 $\phi$ 换成 $\phi-\psi$。
\end{itemize}

把定理 5.4.4 与定理 5.4.5 结合起来，我们立即得到

\begin{theorem}{5.4.6}
[51, Thm. 1.4]（统一 $L^\infty$-估计）给定 Calabi--Yau 流形的大复结构极限，则相对于固定 Fubini--Study 参考度量 $\omega_{FS,t}$、属于类 $\frac{1}{|\log|t||}c_1(L)$ 的 Calabi--Yau 度量 $\omega_{CY,t}$ 的势，在适当的加性归一化之下具有与 $0<|t|\ll1$ 无关的一致 $L^\infty$-估计。
\end{theorem}

我们对 $L^1$-体积稳定性估计的改造成果是

\begin{theorem}{5.4.7}
（参看 [53, Theorem 2.6]）（统一 $L^1$-稳定性）设 $(Y,\omega)$ 是紧 Kähler 流形，$\phi\in PSH(Y,\omega)\cap C^0$，并满足复 MA 方程
\[
\frac{\omega^n}{\mathrm{Vol}(Y)}=d\mu,\qquad \frac{\omega_\phi^n}{\mathrm{Vol}(Y)}=d\nu
\]
其中 $d\mu$、$d\nu$ 是概率测度。假设
\begin{itemize}
\item 存在 Skoda 型估计
\[
\int_Y e^{-\alpha u}d\mu\le A,\qquad \forall u\in PSH(Y,\omega)\ \text{且}\ \sup_Y u=0.
\]
\item 集合 $E_0=\{\phi>0\}$ 的补集具有质量下界
\[
\int_{E_0^c}d\mu\ge\lambda>0.
\]
\item （$L^1$-稳定性假设）全变差 $\int_Y|d\mu-d\nu|\le s^{2n+3}<1$。
\item $\phi$ 在一个（可能为空的）闭子集 $S$（$d\mu$-测度为零）之外光滑，且整体有 $\|\phi\|_{C^0}\le A'$。
\end{itemize}
\end{theorem}

% [原文 p.22]
则对 $0<s<s_0(\lambda,n,\alpha,A,A')\ll1$，存在一致估计
\[
\sup_Y\phi\le C(\lambda,n,\alpha,A,A')s.
\]

定理 5.4.7 应看作 $L^1$-体积稳定性估计的单侧版本。目标是不失一般性地比较两个势 $\phi$ 与 $0$。Skoda 型估计一如前述是体积密度可积性假设 (5.10) 的弱化版本。暂时假定它对两个测度 $d\mu$ 与 $d\nu$ 都成立。在把 $\phi$ 调整一个加性常数之后，我们不妨假设 $\mu(E_0)\approx\nu(E_0)\approx\frac12$，注意由假设两个测度之间的全变差很小。于是我们可以把 $\omega_\phi$ 与 $\omega$ 的角色互换，从而推出关于 $\phi$ 的双侧小性界，这正是 $L^1$-体积稳定性估计的内容。

% [Tutor Note: 定理 5.4.7 的假设中，$E_0^c=\{\phi\le0\}$ 是 $E_0=\{\phi>0\}$ 的补集（原文写 $E_0^c$，注意不是 $E_0$）；$\mu(E_0)\approx\nu(E_0)\approx\frac12$ 只是论证中的“不妨假设”，不是定理的假设。]

\subsection{Savin 小扰动定理}

Savin [65] 证明：对满足某些结构条件的一大类二阶椭圆方程，任何与给定光滑解 $C^0$-接近的粘性解都具有内部 $C^{2,\gamma}$-界。特别地，这适用于复 MA 方程。结合 Schauder 估计，有

\begin{theorem}{5.4.8}
固定 $k\ge2$ 与 $0<\gamma<1$。在单位球上，设 $v$ 是复 Monge--Ampère 方程 $(dd^cv)^n=1$ 的一个给定光滑解。则存在常数 $0<\epsilon\ll1$ 与 $C$（依赖于 $n,k,\gamma,\|v\|_{C^{k,\gamma}}$），使得若
\[
(dd^c(u+v))^n=1+f,\qquad \|f\|_{C^{k-2,\gamma}}<\epsilon,
\]
且 $\|u\|_{C^0}<\epsilon$，则 $\|u\|_{C^{k,\gamma}(B_{1/2})}\le C\epsilon$。
\end{theorem}

% [Tutor Note: 定理 5.4.8 中 $f$ 的正则性是 $C^{k-2,\gamma}$（比 $u$ 所获的 $C^{k,\gamma}$ 低两阶、并少一个 $\gamma$ 的“缺口”），而 $u$ 只需 $C^0$ 小；$v$ 则是给定的光滑解。抄写时不要把 $k-2$ 写成 $k-1$。]

Savin 的定理具有完全非线性性质，因为这一扰动机制只有在解已有先验 $C^2$-界之后才能运用。他的证明有两个主要部分：首先用 Aleksandrov--Bakelman--Pucci 估计作一个非平凡的应用来证明 Harnack 不等式；然后使用紧性论证来证明 $C^{2,\gamma}$-估计，这类似于 De Giorgi 关于极小曲面的几乎平坦性定理 [21]。

\subsection{实 Monge--Ampère 的正则性理论}

关于实 Monge--Ampère 方程局部正则性理论的文献非常丰富，主要归功于 Caffarelli 学派。这里综述的所有结果都可以在 [61] 中找到。

开集 $\Omega\subset\R^n$ 上的任意凸函数 $v:\Omega\to\R$ 都伴随一个称为 Monge--Ampère 测度的 Borel 测度，定义为
\[
MA(v)(E)=|\partial v(E)|,
\]

% [原文 p.23]
其中 $|\partial v(E)|$ 表示次梯度映射在 $E\subset\Omega$ 上像的 Lebesgue 测度。给定 Borel 测度 $\mu$，$MA(v)=\mu$ 的解称为 $\det(D^2v)=\mu$ 的\emph{Aleksandrov 解}；若 $v\in C^2$，这就是经典的实 Monge--Ampère 方程。我们将假定双侧密度界
\[
\det(D^2v)=f\ \text{在}\ B_1\ \text{中},\qquad 0<\Lambda_1\le f\le\Lambda_2.
\]
设 $B_1\setminus\Sigma$ 是 $v$ 的严格凸点集，即存在一个支撑超平面只在一个点处与 $v$ 的图像相触。于是 Caffarelli [7][8][9] 证明
\begin{itemize}
\item 若 $f\in C^\gamma(B_1)$，则 $v\in C^{2,\gamma}_{loc}(B_1\setminus\Sigma)$。再由 Schauder 理论，若 $f$ 光滑，则 $v$ 在 $B_1\setminus\Sigma$ 中光滑。
\item 若 $L$ 是 $v$ 的一个支撑仿射线性函数，使得凸集 $\{v=L\}$ 不是一个点，则 $\{v=L\}$ 在 $B_1$ 的内部没有极点。
\item 上述仿射线性集 $\{v=L\}$ 的维数 $k<n/2$。
\end{itemize}
Mooney [61] 进一步证明
\begin{itemize}
\item 奇异集 $\Sigma$ 具有 $(n-1)$-Hausdorff 测度零。从而 $B_1\setminus\Sigma$ 是道路连通的（因为连接两个给定点的一般道路不会与一个零 $(n-1)$-Hausdorff 测度的子集相交）。
\item 即使 $\Sigma$ 非空，解也有 $v\in W^{2,1}_{loc}(B_1)$。
\end{itemize}

\begin{remark}{5.8}
Pogorelov 的一个经典反例表明：当 $n=3$ 时，奇异集 $\Sigma$ 可以包含一条线段。Caffarelli [9] 对此作了推广：对任意 $k<n/2$，他构造了 $f$ 光滑但 $\Sigma$ 包含一个 $k$-平面的例子。Mooney [61] 的一个出人意料的例子表明，$\Sigma$ 的 Hausdorff 维数可以大于 $n-1-\epsilon$（对任意小的 $\epsilon$）。这意味着上面综述的局部正则性理论本质上已是最优的。
\end{remark}

% [Tutor Note: 注意“严格凸点”的定义用的是“支撑超平面只在一点接触图像”；$\Sigma$ 是奇异集，Caffarelli 的三条结论分别说明：$\Sigma$ 之外的高阶正则性、$\{v=L\}$ 无内点极点、以及仿射线集的维数 $k<n/2$。Remark 5.8 中 $n-1-\epsilon$ 的“几乎 $n-1$ 维”正说明前面“$(n-1)$-Hausdorff 测度零”的估计不能改进太多。]

\subsection{特殊拉格朗日纤维化}

McLean 的经典结果表明：相位为 $\theta$ 的特殊拉格朗日的形变理论是无障碍的，且一阶形变空间同构于 $H^1(L,\R)$。因此若 $L$ 微分同胚于 $T^n$，则形变空间是 $n$ 维的，这与 SYZ 猜想中 $X$ 容许一个特殊拉格朗日 $T^n$-纤维化相符。在度量坍缩且截面曲率局部有界这一非常强的假设下，Zhang [80, Thm 1.1] 得到了构造 Slag 纤维化的一个充分条件。

Zhang 结果的实质是隐函数定理的标准应用，我们将概述其中的要点（细节参看 [80, 第 4 节]）。

% [Tutor Note: “Slag” 是 special Lagrangian 的常用缩写（原文 p.23 首次使用 Slag 一词）；译文按 glossary 统一作“特殊拉格朗日”，仅在本段保留缩写形式以便与原文对照。]
